% Visueller Regressionstest gegen 1127501_WB00-BAC-PB1-260630.pdf.
%
% Diese Datei kopiert die relevanten Formatierungsdefinitionen aus ../main.tex.
% Sie dient nur dem Vergleich und darf keine produktiven Vorgaben veraendern.
\documentclass[
  12pt,
  a4paper,
  oneside,
  headings=normal,
  numbers=noendperiod
]{scrreprt}

\newcommand{\MainLanguage}{ngerman}
\newcommand{\ThesisType}{Bachelorarbeit}
\newcommand{\ThesisTitle}{Visueller Layoutvergleich}
\newcommand{\AuthorName}{Marc-William Wagner}
\newcommand{\MatriculationNumber}{1127501}

\usepackage[T1]{fontenc}
\usepackage[utf8]{inputenc}
\usepackage[english,ngerman]{babel}
\usepackage{newtxtext,newtxmath}
\usepackage{xcolor}
\usepackage{csquotes}
\usepackage{etoolbox}

\usepackage[
  a4paper,
  portrait,
  left=3cm,
  right=4cm,
  top=2.5cm,
  bottom=2.5cm
]{geometry}

\usepackage{setspace}
\setstretch{1.43}
\setlength{\parindent}{0pt}
\setlength{\parskip}{6pt}

\AtBeginDocument{\selectlanguage{\MainLanguage}}
\newcommand{\DEEN}[2]{\iflanguage{english}{#2}{#1}}
\setcounter{secnumdepth}{2}
\setcounter{tocdepth}{2}

\setkomafont{disposition}{\normalfont\bfseries\setstretch{1}}
\setkomafont{chapter}{\fontsize{14}{17}\selectfont}
\setkomafont{section}{\fontsize{12}{14.5}\selectfont}
\setkomafont{subsection}{\fontsize{12}{14.5}\selectfont}
\setkomafont{subsubsection}{\fontsize{12}{14.5}\selectfont}
\RedeclareSectionCommand[beforeskip=0pt,afterskip=18pt]{chapter}
\RedeclareSectionCommand[beforeskip=18pt,afterskip=12pt]{section}
\RedeclareSectionCommand[beforeskip=18pt,afterskip=12pt]{subsection}
\RedeclareSectionCommand[beforeskip=18pt,afterskip=12pt]{subsubsection}

\usepackage[automark,singlespacing=true]{scrlayer-scrpage}
\clearpairofpagestyles
\ihead{\AuthorName}
\ohead{\MatriculationNumber}
\cfoot{}
\AtBeginDocument{%
  \cfoot{Seite \thepage\ von \pageref*{LastNumberedPage}}}
\setkomafont{pageheadfoot}{%
  \normalfont\fontsize{10}{12}\selectfont}
\setlength{\headheight}{15pt}
\pagestyle{scrheadings}
\renewcommand*{\chapterpagestyle}{scrheadings}

\usepackage{chngcntr}
\counterwithout{footnote}{chapter}
\setlength{\footnotesep}{16pt}
\newcommand*{\HFHFootnoteMark}{%
  \raisebox{0.7ex}{\fontsize{9}{9}\selectfont\thefootnotemark}}
\deffootnotemark{\HFHFootnoteMark}
\deffootnote[0.5cm]{0.5cm}{0cm}{%
  \raggedright\makebox[0.5cm][l]{\HFHFootnoteMark}}

\usepackage{graphicx}
\usepackage{booktabs}
\usepackage{tabularx}
\usepackage{array}
\usepackage{caption}
\captionsetup{
  format=plain,
  justification=justified,
  singlelinecheck=false,
  font={footnotesize,it},
  labelfont={bf,it},
  labelsep=colon,
  skip=6pt
}
\counterwithout{figure}{chapter}
\counterwithout{table}{chapter}
\setlength{\intextsep}{12pt}
\setlength{\textfloatsep}{12pt}
\numberwithin{equation}{chapter}

\newenvironment{hfhtable}[1][htbp]
  {\begin{table}[#1]\centering\footnotesize\setstretch{1}}
  {\end{table}}
\newenvironment{hfhfigure}[1][htbp]
  {\begin{figure}[#1]\centering\footnotesize\setstretch{1}}
  {\end{figure}}
\newcommand{\source}[1]{%
  \par\vspace{3pt}{\footnotesize\itshape\setstretch{1}%
    \DEEN{Quelle}{Source}: #1\par}}
\newcolumntype{Y}{>{\raggedright\arraybackslash}X}

\usepackage[round,authoryear]{natbib}
\setlength{\bibhang}{0.5cm}
\setlength{\bibsep}{6pt}
\AtBeginEnvironment{thebibliography}{\singlespacing}

\newenvironment{hfhsourceentries}
  {\par\begingroup\singlespacing\setlength{\parskip}{6pt}}
  {\par\endgroup}
\newcommand{\hfhsourceentry}[1]{%
  \par\noindent\hangindent=0.5cm\hangafter=1 #1\par}

\usepackage{xurl}
\usepackage[hidelinks]{hyperref}
\usepackage{bookmark}
\usepackage{pdfpages}
\hypersetup{
  pdftitle={\ThesisTitle},
  pdfauthor={\AuthorName},
  pdfsubject={\ThesisType}
}

\newcommand{\TemplateHint}[1]{%
  \par\begingroup\small\itshape\color{gray}#1\par\endgroup}

\begin{document}

% Vergleich A: Originalseiten 6-7.
\phantomsection
\label{test:chapter-start}
\chapter{Einleitung}

In modernen Industrieunternehmen spielt ein strukturiertes und sicheres
Zutrittsmanagement eine zentrale Rolle für den Schutz sensibler
Unternehmensbereiche sowie für die Einhaltung regulatorischer Anforderungen.
Insbesondere in der pharmazeutischen Industrie unterliegen Produktionsstandorte
strengen gesetzlichen und organisatorischen Vorgaben, die sowohl den Zugang zu
Produktionsanlagen als auch zu Labor- und Lagerbereichen betreffen. Neben
Sicherheitsaspekten müssen dabei auch Anforderungen an Dokumentation,
Nachvollziehbarkeit sowie Datenschutz berücksichtigt werden.

\section{Problemstellung und Zielsetzung}

Am Pharma-Produktionsstandort Wuppertal der Bayer AG erfolgt die Vergabe von
Zutrittsberechtigungen für interne Mitarbeitende, externe Dienstleister sowie
Besucher derzeit überwiegend über manuelle und papierbasierte Prozesse.
Angesichts einer großen Anzahl von etwa 4200 Personen (Interne und Externe
Mitarbeitende), die regelmäßig Zugang zum Werksgelände benötigen, führt diese
Vorgehensweise zu einem höheren administrativen Aufwand und zu langen
Bearbeitungszeiten im Genehmigungsprozess. Gleichzeitig erschwert die manuelle
Prozessgestaltung eine transparente Nachvollziehbarkeit von Entscheidungen sowie
eine effiziente Dokumentation der vergebenen Zutrittsrechte.

Besonders relevant ist diese Problematik vor dem Hintergrund regulatorischer
Anforderungen im pharmazeutischen Produktionsumfeld. Produktionsstandorte
unterliegen den Vorgaben der Good Manufacturing Practice (GMP), welche eine
lückenlose Dokumentation und Kontrolle aller relevanten Prozesse innerhalb der
Produktionsumgebung verlangen. Das Unternehmen Bayer produziert für den
internationalen Markt und muss neben den europäischen GMP-Vorgaben auch
internationale regulatorischer Anforderungen berücksichtigen. Dazu zählen die
Regularien der US-amerikanischen Food and Drug Administration (FDA) sowie die
Richtlinien der International Council für Harmonisation (ICH). Da über
Zutrittsberechtigungen der Zugang zu GMP-relevanten Bereichen gesteuert wird,
muss auch der zugehörige Genehmigungsprozess entsprechend strukturiert,
nachvollziehbar und auditfähig gestaltet sein.

Aufgrund der am Standort verarbeiteten Gefahrstoffe wird dieser in der hohen
Störfallkategorie eingestuft, wodurch ein erhöhtes Sicherheitsniveau und eine
differenzierte und streng kontrollierte Verwaltung von Zutrittsrechten
erforderlich wird. Zusätzlich wurde der Standort unternehmensintern als BCC4
(Business Criticality Classification) eingestuft, diese Klassifizierungen und
die hohe Störfallkategorie werden im weiteren Verlauf näher erklärt.

Die derzeitige Prozessgestaltung führt unter anderem zu eingeschränkter
Prozessübersicht, erhöhter Fehleranfälligkeit sowie zu Herausforderungen bei der
Daten\-schutz- Grundverordnung (DSGVO) konformen Dokumentation personenbezogener
Daten. Gleichzeitig steigen mit der fortschreitenden Digitalisierung
betrieblicher Abläufe die Erwartungen an effiziente, transparente und
softwaregestützte Prozesslösungen. Vor diesem Hintergrund verfolgt die
vorliegende Bachelorarbeit das Ziel, den bestehenden Prozess der
Zutrittsvergabe systematisch zu analysieren und auf dieser Grundlage ein Konzept
für eine digitale und effizientere Prozessgestaltung zu entwickeln. Dabei sollen
Anforderungen an eine geeignete Softwarelösung identifiziert, potenzielle
Systeme bewertet sowie ein digitaler Soll-Prozess für die Vergabe von
Zutrittsberechtigungen konzipiert werden. Die bestehende physische
Zutrittsinfrastrukturen wie Zugangstore, Kartenleser oder Schließsysteme bleiben
davon unverändert bestehen und werden lediglich durch digitale
Verwaltungsprozesse ergänzt.

Ein weiteres Ziel der Arbeit besteht darin, ein Konzept zur Einführung eines
entsprechenden digitalen Systems zu entwickeln, das sowohl organisatorische als
auch technische Aspekte berücksichtigt. Dadurch soll eine Grundlage geschaffen
werden, auf deren Basis der Zutrittsprozess langfristig effizienter, transparenter
und revisionssicher gestaltet werden kann.

\section{Motivation und Relevanz des Themas}

Die Digitalisierung betrieblicher Prozesse stellt einen zentralen Bestandteil
moderner Unternehmensentwicklung dar. Durch den Einsatz digitaler Technologien
können Arbeitsabläufe effizienter gestaltet, Fehlerquellen reduziert und
Transparenz innerhalb organisatorischer Prozesse erhöht werden. Insbesondere in
sicherheitsrelevanten Bereichen wie dem Zutrittsmanagement kommt der
Digitalisierung eine besondere Bedeutung zu, da hier neben Effizienzsteigerungen
auch Aspekte der Unternehmenssicherheit und Compliance berücksichtigt werden
müssen.

\clearpage

% Vergleich B: Originalseiten 18-19.
\phantomsection
\label{test:pure-prose}
technische Einrichtungen manipuliert werden. Zugangskontrollen,
Sicherungssysteme und räumliche Abschirmungen sind daher geeignet, zur Erfüllung
der gesetzlichen Betreiberpflichten beizutragen. Im Rahmen der
Störfallverordnung wird zwischen Betriebsbereichen unterschiedlicher
Gefährdungsklassen unterschieden. Ein Betriebsbereich liegt vor, wenn
gefährliche Stoffe in relevanten Mengen vorhanden sind oder entstehen können.
Dabei wird zwischen der unteren und der oberen Störfallklasse differenziert.
Betriebsbereiche der oberen Klasse zeichnen sich durch ein erhöhtes
Gefährdungspotenzial aus, da größere Mengen gefährlicher Stoffe vorhanden sind.
Daraus resultieren strengere Anforderungen an Sicherheitsmaßnahmen, Überwachung
und organisatorische Strukturen (vgl. BImSchG 2025: § 3 Abs. 5a).

Im Kontext kritischer Infrastrukturen kommt den Betreibern eine zentrale
Verantwortung für die Gewährleistung der Sicherheit und Funktionsfähigkeit ihrer
Anlagen zu. Das KRITIS-Dachgesetz definiert Betreiber kritischer Anlagen, die
einen maßgeblichen Einfluss auf Einrichtungen ausüben, deren Ausfall zu
erheblichen Versorgungsengpässen oder Gefährdungen der öffentlichen Sicherheit
führen kann. Aus dieser besonderen Stellung resultiert eine umfassende
Verpflichtung zur Sicherstellung der Resilienz kritischer Anlagen. Resilienz wird
dabei als Fähigkeit verstanden, Vorfälle zu verhindern, sich vor ihnen zu
schützen, auf sie zu reagieren, deren Auswirkungen zu begrenzen und eine
schnelle Wiederherstellung der Funktionsfähigkeit zu gewährleisten. Betreiber
sind somit nicht nur für den sicheren Betrieb verantwortlich, sondern auch für
präventive Gestaltung robuster und widerstandsfähiger Systeme. (vgl.
KRITISDachG 2026: 2 f.).

Der Schutz vor unbefugten Eingriffen stellt eine Anforderung im
sicherheitskritischen Anlagenbetrieb dar und ist im Kontext der
Störfall-Verordnung verpflichtend zu berücksichtigen. Betreiber von Anlagen
sind dazu verpflichtet, geeignete Maßnahmen zu ergreifen, um Gefahren durch
vorsätzliche Eingriffe Dritter zu verhindern und potenzielle Schadensereignisse
zu minimieren. Ziel ist es, sicherzustellen, dass gefährliche Stoffe so
gesichert sind, dass ernsthafte Gefahren für Mensch und Umwelt weitestgehend
ausgeschlossen werden können (vgl. KAS 2019: 2 f.). Dabei umfasst der Begriff
der unbefugten Eingriffe sowohl physischen Zutritt als auch den Zugriff auf
Systeme und Daten. Entsprechend müssen Unternehmen organisatorische als auch
technische Maßnahmen implementieren, die den Zugang zu sensiblen Bereichen
kontrollieren und überwachen. Das Zutritts- und Zugriffsmanagement spielt
hierbei eine zentrale Rolle, da physische Zugänge zu Anlagen und digitale
Zugriffe auf IT- und Steuerungssysteme gezielt eingeschränkt und dokumentiert
werden müssen (vgl. KAS 2019: 9 f.).

Im Rahmen des unternehmensweiten Sicherheitsmanagements der Bayer AG dient die
Business Criticality Classification als zentrales Instrument zur Priorisierung
von Standorten. Als interne Richtlinie der Corporate Global Security ermöglicht
sie eine systematische Bewertung der geschäftlichen Bedeutung einzelner
Standorte, um Sicherheitsmaßnahmen und Ressourcen risikoorientiert auszurichten.
Die BCC unterscheidet vier Kritikalitätsstufen (BCC 1 bis BCC 4), die sich nach
der Bedeutung eines Standorts für die Geschäftstätigkeit richten. Standorte mit
hoher Kritikalität unterliegen dabei häufigeren Überprüfungen als weniger
kritische Standorte, wodurch eine effiziente Allokation von
Sicherheitsressourcen gewährleistet wird. Die Einstufung erfolgt im Rahmen eines
standardisierten Bewertungsprozesses unter Beteiligung von Risk Ownern sowie
Sicherheitsverantwortlichen. Grundlage bildet ein strukturierter
Selbstbewertungsfragebogen, dessen Ergebnisse zentral dokumentiert werden. Die
BCC stellt somit ein Steuerungsinstrument dar, das eine nachvollziehbare
Priorisierung von Standorten ermöglicht und zu gezielten Absicherung
geschäftskritischer Prozesse innerhalb von Bayer beiträgt (vgl. Grynchak 2024:
8 f.).

Im Kontext der Unternehmenssicherheit wird der Schutz von Gebäuden, Anlagen und
Informationen häufig anhand eines mehrschichtigen Sicherheitsmodells
beschrieben. Dieses sogenannte Schalenmodell unterteilt Sicherheitsmaßnahmen in
verschiedene Ebenen, die den Zugang zu sicherheitskritischen Bereichen
schrittweise kontrollieren. Eine zentrale Rolle spielt dabei der Zutrittsschutz,
der den physischen Zugang zu Gebäuden und sensiblen Bereichen regelt. Durch
technische Systeme wie Zutrittskontrollanlagen wird sichergestellt, dass
ausschließlich autorisierte Personen Zugang erhalten. Ergänzend dazu umfasst der
Zugangsschutz die Absicherung von IT-Systemen durch
Authentifizierungsmechanismen und technischen Schutzmaßnahmen. Der
Zugriffsschutz stellt sicher, dass nur berechtigte Personen auf bestimmte Daten
und Anwendungen zugreifen können (vgl. Müller 2023: 418).

\clearpage

% Vergleich C: Originalseiten 20-21.
\phantomsection
\label{test:subsections}
\setcounter{chapter}{2}
\setcounter{section}{3}
\setcounter{subsection}{1}
\subsection{Schutz vor unbefugtem Zutritt}

Zutrittskontrollsysteme stellen einen wesentlichen Bestandteil der betrieblichen
Sicherheitsinfrastruktur dar. Sie dienen dazu, den Zugang zu Gebäuden, Anlagen
und sensiblen Bereichen gezielt zu steuern und unbefugten Zutritt zu verhindern.
Ein unkontrollierter Zugang stellt ein erhebliches Risiko für Unternehmen dar, da
er potenzielle Gefahren wie Diebstahl, Sabotage oder Wirtschaftsspionage
begünstigt. Im Kern basiert ein Zutrittskontrollsystem auf dem Prinzip der
gezielten Steuerung von Personenströmen. Dabei wird festgelegt, welche Personen
Zugang zu bestimmten Bereichen erhalten. Dieses Steuerungsprinzip lässt sich
durch die drei zentralen Kriterien „Wer“, „Wann“ und „Wo“ beschreiben.
Zutrittskontrollsysteme ermöglichen eine Verwaltung von Zutrittsrechten, diese
können differenziert nach Zeit- und Raumprofilen vergeben werden (vgl. BHE
Bundesverband Sicherheitstechnik e.V. 2025: 1).

Ein wesentlicher Bestandteil eines wirksamen Zutrittsmanagements ist die
Sicherstellung der Aktualität sowie der vollständigen Dokumentation von
Zutrittsberechtigungen. Dabei kommt insbesondere der Pflege von
Zutrittsregelungen und zugehörigen Berechtigungslisten eine zentrale Bedeutung
zu. Diese müssen jederzeit aktuell gehalten werden, um sicherzustellen, dass
ausschließlich berechtigte Personen Zugang zu sicherheitsrelevanten Bereichen
erhalten (vgl. Schwarz: 164).

\subsection{Informationssicherheit und Identitäts- und Zugriffsmanagement}

Die Informationssicherheit stellt eine zentrale Voraussetzung für den sicheren
Betrieb moderner Organisationen dar, da nahezu alle Geschäftsprozesse von
informationstechnischen Systemen abhängig sind. Ziel der Informationssicherheit
ist es, Informationen vor unbefugtem Zugriff, Manipulation oder Ausfall zu
schützen. Dabei stehen die Schutzziele Vertraulichkeit, Integrität und
Verfügbarkeit im Fokus, deren Verletzung zu erheblichen betrieblichen und
wirtschaftlichen Schäden führen kann (vgl. BSI 2023: 1).

Im Zuge von Industrie 4.0 werden Produktionsanlagen, IT-Systeme und digitale
Geschäftsprozesse zunehmend miteinander vernetzt. Dadurch entstehen neue
Anforderungen an die IT-Sicherheit, da Angriffe auf vernetzte Systeme nicht nur
digitale Daten, sondern auch physische Anlagen und Produktionsprozesse
beeinträchtigen können. In den industriellen Umgebungen müssen Systeme vor
unbefugten Zugriffen, Manipulationen und Datenverlust geschützt werden.
Vernetzte Komponenten und Systeme benötigen daher geeignete
Sicherheitsmechanismen zur Authentifizierung, Autorisierung und sicheren
Kommunikation. Eine zentrale Rolle spielt hierbei das Berechtigungsmanagement,
durch das festgelegt wird, welche Personen oder Systeme auf bestimmte Funktionen
und Daten zugreifen dürfen. Darüber hinaus müssen Sicherheitsmaßnahmen gegen
digitale als auch gegen physische Angriffe ausgelegt sein, um die Integrität und
Verfügbarkeit der Systeme sicherzustellen (vgl. Vogel -- Heuser et al 2016:
135 f.).

Das Identitäts- und Berechtigungsmanagement (Identity and Access Management --
IAM) stellt einen Bestandteil der Informationssicherheit dar und dient der
kontrollierten Vergabe und Verwaltung von Zugriffsrechten innerhalb einer
Organisation. Es stellt sicher, dass autorisierte Personen Zugriff auf
Informationen, Systeme und physische Ressourcen erhalten und somit unbefugte
Nutzung verhindert wird. Im IT-Grundschutz-Kompendium des Bundesamts für
Sicherheit in der Informationstechnik (BSI) wird gefordert, dass Berechtigungen
strukturiert und nachvollziehbar vergeben werden. Hierbei ist sicherzustellen,
dass Identitäten eindeutig definiert und den jeweiligen Rollen und
Verantwortlichkeiten zugeordnet werden. Dies ermöglicht eine klare Abgrenzung
von Zuständigkeiten und reduziert das Risiko von Fehlkonfigurationen oder
missbräuchlicher Nutzung. Darüber hinaus wird verlangt, dass Prozesse zur
Einrichtung, Änderung und Löschung von Benutzerkonten definiert und umgesetzt
werden. Diese Prozesse müssen sicherstellen, dass Berechtigungen jederzeit
aktuell sind und bei Rollenwechseln oder dem Ausscheiden von Mitarbeitenden
zeitnah angepasst oder entzogen werden. Eine regelmäßige Überprüfung bestehender
Berechtigungen ist hierbei unerlässlich, um übermäßige oder nicht mehr benötigte
Zugriffsrechte zu identifizieren und zu entfernen (vgl. BSI 2023: ORP.4).

Neben hierarchischen und netzwerkbasierten Datenbankmodellen stellen relationale
Datenbanken heute die Grundlage vieler moderner Informationssysteme dar. Das
relationale Datenbankmodell wurde 1970 von Edgar Codd entwickelt und basiert auf
der strukturierten Speicherung von Daten in Tabellen mit Zeilen und Spalten
(vgl. Taylor, Blum 2025: 50). Im Gegensatz zu hierarchischen Datenbanken, die
Daten in festen Baumstrukturen organisieren, und Netzwerkdatenbanken, die
komplexe Verknüpfungen zwischen Datensätzen ermöglichen, zeichnen sich
relationale Datenbanken durch eine flexible und strukturierte Verwaltung von
Daten aus.

\clearpage

% Vergleich D: Originalseiten 34-35.
\phantomsection
\label{test:nested-chapter}
\setcounter{chapter}{3}
\chapter{Systemanalyse und Soll-Konzeption}

Auf Grundlage der in Kapitel 3 identifizierten Schwachstellen des bestehenden
Zutrittsvergabeprozesses sowie der daraus abgeleiteten Anforderungen werden in
diesem Kapitel untersucht, welche digitalen Lösungen zur Unterstützung des
Zutrittsmanagements geeignet sind. Ziel ist es, Systeme zu identifizieren, die
eine effiziente, transparente und nachvollziehbare Verwaltung von
Zutrittsberechtigungen ermöglichen und gleichzeitig die organisatorischen,
technischen sowie regulatorischen Anforderungen des Standortes erfüllen.

Neben der Analyse möglicher Systemlösungen wird im weiteren Verlauf des
Kapitels ein digitaler Soll-Prozess entwickelt. Dabei erfolgt die Systemanalyse
zunächst auf Basis der identifizierten Anforderungen, um geeignete technische
Lösungsansätze zu bestimmen. Aufbauend auf diesen Ergebnissen wird anschließend
ein Soll-Prozess konzipiert, der die funktionalen Möglichkeiten der Systeme als
auch die spezifischen Anforderungen des Standorts berücksichtigt.

\section{Bewertung möglicher Systeme}

Im Unternehmensumfeld lassen sich unterschiedliche Systemtypen identifizieren,
die für die Umsetzung eines digitalen Zutrittsmanagements geeignet sein können.
Dazu zählen spezialisierte Zutrittsmanagementsysteme,
Workflow-Management-Systeme und Identity- und Access-Management Lösungen.

\subsection{Identifikation potenzieller digitaler
Zutrittsmanagementlösungen}

Zur Digitalisierung des bestehenden Zutrittsmanagementprozesses kommen
grundsätzlich verschiedene Systemansätze in Betracht, die unterschiedliche
funktionale Schwerpunkte adressieren.

Spezialisierte Zutrittsmanagementsysteme fokussieren sich primär auf die
Verwaltung physischer Zutrittsberechtigungen sowie die Integration mit
technischen Komponenten wie Kartenlesern oder elektronischen Schließsystemen.
Sie ermöglichen die Definition von Zutrittsprofilen und die Steuerung und
Durchsetzung von Zugangsrechten innerhalb physischer Infrastrukturen. Der
Schwerpunkt dieser Systeme liegt somit auf der operativen Umsetzung von
Zutrittsberechtigungen.

Eine weitere Möglichkeit zur Digitalisierung besteht im Einsatz von
Workflow-Management-Systemen. Diese Systeme dienen der strukturierten Abbildung
organisatorischer Prozesse und werden insbesondere zur digitalen Umsetzung von
Genehmigungs- und Freigabeprozessen eingesetzt. Sie ermöglichen die
modellbasierte Definition von Abläufen, die automatisierte Weiterleitung von
Aufgaben zwischen verschiedenen Rollen und die transparente Nachverfolgung von
Prozessschritten.

Darüber hinaus stellen Identity- und Access-Management-Systeme (IAM) eine
zentrale Komponente moderner IT-Landschaften dar. Sie dienen der Verwaltung
digitaler Identitäten sowie der Zuordnung von Rollen und Berechtigungen. Durch
den Einsatz rollenbasierter Zugriffskonzepte können Nutzergruppen und
Verantwortlichkeiten systematisch abgebildet und verwaltet werden. Der Fokus
dieser Systeme liegt auf der Governance von Identitäten und digitalen
Zugriffsrechten.

\subsection{Ableitung eines geeigneten Systemansatzes}

Ausgehend von den Anforderungen an das digitale Zutrittsmanagement ergeben sich
zentrale funktionale und organisatorische Kriterien. Dazu zählen die
strukturierte Abbildung von Genehmigungsprozessen, die Einbindung mehrere
organisatorischen Rollen, die Nachvollziehbarkeit von Entscheidungen und die
revisionssichere Dokumentation aller Prozessschritte. Darüber hinaus sind die
Integration in bestehende Systeme und die Einhaltung regulatorischer
Anforderungen von Bedeutung.

Am Standort Wuppertal ist bereits ein Zutrittskontrollsystem der Firma dormakaba
im Einsatz, dass die technische Umsetzung und Steuerung physischer
Zutrittsberechtigungen übernimmt. Dieses System bildet die operative Grundlage
für die Vergabe und Durchsetzung von Zutrittsrechten. Allerdings ist das
bestehende Zutrittskontrollsystem primär auf die technische Verwaltung und
Ausführung von Zutrittsrechten ausgelegt und bietet keine umfassende
Unterstützung für einen strukturierten digitalen
Berechtigungsvergabeprozess. Es fehlen Funktionen zur standardisierten
Beantragung, workflowbasierten Genehmigung, revisionssicheren Dokumentation
sowie zur Abbildung komplexer Rollen- und Verantwortlichkeitsstrukturen. Darüber
hinaus können organisatorische Anforderungen, wie die transparente
Nachverfolgung von Genehmigungsprozessen, nur eingeschränkt umgesetzt werden.

\phantomsection
\label{LastNumberedPage}

\end{document}
